\chapter{Relazioni d'ordine}\label{cap:ordini}

Le relazioni d'ordine formalizzano l'idea di «venire prima» o di «essere più piccolo»: $\leq$ tra i numeri, $\subseteq$ tra
gli insiemi, «divide» tra i numeri naturali. A differenza dei numeri, però, in un ordine generico due elementi possono anche
\textbf{non essere confrontabili}: $\{1\}$ e $\{2\}$ non sono uno contenuto nell'altro. In questo capitolo vediamo il legame tra
ordini larghi e stretti, gli elementi «estremi» (minimo, massimo, minimali, massimali), le funzioni che rispettano l'ordine e
un modo comodo di disegnare gli insiemi ordinati finiti: i \textbf{diagrammi di Hasse}.

% ──────────────────────────────────────────────────────────────
\section{Ordini larghi e ordini stretti}\label{sec:largo-stretto}

Ricordiamo (Definizione~\ref{def:tipi-relazioni}) che un ordine \textbf{largo} è riflessivo, antisimmetrico e transitivo, come
$\leq$, mentre un ordine \textbf{stretto} è antiriflessivo e transitivo, come $<$. Da un tipo si passa all'altro aggiungendo o
togliendo l'uguaglianza.

\dfn[ol-os]{Gli insiemi $\mathrm{OL}(s)$ e $\mathrm{OS}(s)$}{
    Sia $s$ un insieme. Indichiamo con $\mathrm{OL}(s)$ l'insieme delle relazioni d'ordine largo su $s$ e con $\mathrm{OS}(s)$
    quello delle relazioni d'ordine stretto.
    \begin{itemize}
        \item Se $\rho \in \mathrm{OL}(s)$, definiamo la relazione $\rho^\wedge$ su $s$ ponendo \ $x \rho^\wedge y \leftrightarrow (x \rho y \wedge x \neq y)$.
        \item Se $\rho \in \mathrm{OS}(s)$, definiamo la relazione $\rho^\vee$ su $s$ ponendo \ $x \rho^\vee y \leftrightarrow (x \rho y \vee x = y)$.
    \end{itemize}
}

\begin{concetto}
\begin{Teorema}{Ordini larghi e ordini stretti si corrispondono}{largo-stretto}
    \begin{prerequisiti}
        \dip{Definizione}{def:ol-os}
        \dip{Definizione}{def:tipi-relazioni}
        \dip{Proposizione}{prop:distributivita}
        \dip{Teorema}{teo:biettiva-inversa}
    \end{prerequisiti}
    Sia $s$ un insieme. Allora $\rho \in \mathrm{OL}(s) \mapsto \rho^\wedge \in \mathrm{OS}(s)$ è una funzione biettiva, e la sua inversa è
    $\rho \in \mathrm{OS}(s) \mapsto \rho^\vee \in \mathrm{OL}(s)$.
\end{Teorema}
\begin{dimostrazione}
    \facoltativa
    \begin{caso}{Se $\rho \in \mathrm{OL}(s)$, allora $\rho^\wedge \in \mathrm{OS}(s)$}
        \textbf{Antiriflessiva}: $x \rho^\wedge x$ richiederebbe $x \neq x$, che è falso. \textbf{Transitiva}: siano $x \rho^\wedge y$ e
        $y \rho^\wedge z$, cioè $x \rho y$, $x \neq y$, $y \rho z$, $y \neq z$. Per la transitività di $\rho$, $x \rho z$. Se fosse $x = z$,
        avremmo $x \rho y$ e $y \rho x$, quindi $x = y$ per l'antisimmetria: assurdo. Dunque $x \neq z$ e $x \rho^\wedge z$.
    \end{caso}
    \begin{caso}{Se $\rho \in \mathrm{OS}(s)$, allora $\rho^\vee \in \mathrm{OL}(s)$}
        \textbf{Riflessiva}: $x = x$. \textbf{Antisimmetrica}: siano $x \rho^\vee y$ e $y \rho^\vee x$; se fosse $x \neq y$ avremmo
        $x \rho y$ e $y \rho x$, quindi $x \rho x$ per transitività, contro l'antiriflessività; dunque $x = y$. \textbf{Transitiva}: siano
        $x \rho^\vee y$ e $y \rho^\vee z$. Se $x = y$ o $y = z$ la tesi è una delle ipotesi; altrimenti $x \rho y$ e $y \rho z$, quindi $x \rho z$.
    \end{caso}
    \begin{caso}{Le due funzioni sono una l'inversa dell'altra}
        Sia $\rho \in \mathrm{OL}(s)$. Usando la distributività di $\vee$ rispetto a $\wedge$:
        \[ x (\rho^\wedge)^\vee y \leftrightarrow (x \rho y \wedge x \neq y) \vee x = y \leftrightarrow (x \rho y \vee x = y) \wedge (x \neq y \vee x = y)
           \leftrightarrow x \rho y \vee x = y \leftrightarrow x \rho y, \]
        dove l'ultimo passaggio usa la riflessività di $\rho$. Sia $\rho \in \mathrm{OS}(s)$. Usando la distributività di $\wedge$ rispetto a $\vee$:
        \[ x (\rho^\vee)^\wedge y \leftrightarrow (x \rho y \vee x = y) \wedge x \neq y \leftrightarrow (x \rho y \wedge x \neq y) \vee (x = y \wedge x \neq y)
           \leftrightarrow x \rho y \wedge x \neq y \leftrightarrow x \rho y, \]
        dove l'ultimo passaggio usa l'antiriflessività di $\rho$ ($x \rho y$ implica $x \neq y$). Due funzioni una inversa dell'altra
        sono biettive (Teorema~\ref{teo:biettiva-inversa}).
    \end{caso}
\end{dimostrazione}
\end{concetto}

\begin{spiegazione}
    Sul grafo con frecce, passare da un ordine largo a uno stretto vuol dire \textbf{togliere tutti i cappi}; tornare indietro vuol
    dire \textbf{rimetterli}. Le altre frecce non cambiano. Ecco $\leq$ e $<$ su $\{1, 2, 3\}$:
    \begin{center}
    \begin{tikzpicture}
        \node[rv] (1) at (0,0) {$1$}; \node[rv] (2) at (1.6,0) {$2$}; \node[rv] (3) at (3.2,0) {$3$};
        \draw[rfreccia] (1) to[rcappio] (1); \draw[rfreccia] (2) to[rcappio] (2); \draw[rfreccia] (3) to[rcappio] (3);
        \draw[rfreccia] (1) to (2); \draw[rfreccia] (2) to (3); \draw[rfreccia] (1) to[bend right=35] (3);
        \node[font=\small] at (1.6,-1.15) {$\leq$ (ordine largo)};
        \draw[-{Stealth}, line width=1pt, primary] (4.0,0.25) -- node[above, font=\small] {$\rho \mapsto \rho^\wedge$} (5.6,0.25);
        \draw[{Stealth}-, line width=1pt, seagreen!65!black] (4.0,-0.25) -- node[below, font=\small] {$\rho^\vee \mapsfrom \rho$} (5.6,-0.25);
        \begin{scope}[xshift=6.4cm]
            \node[rv] (1) at (0,0) {$1$}; \node[rv] (2) at (1.6,0) {$2$}; \node[rv] (3) at (3.2,0) {$3$};
            \draw[rfreccia] (1) to (2); \draw[rfreccia] (2) to (3); \draw[rfreccia] (1) to[bend right=35] (3);
            \node[font=\small] at (1.6,-1.15) {$<$ (ordine stretto)};
        \end{scope}
    \end{tikzpicture}
    \end{center}
    Allo stesso modo $(\subseteq)^\wedge = {\subset}$ e $(\subset)^\vee = {\subseteq}$. Il punto delicato della dimostrazione è che, togliendo i
    cappi, la transitività non si perde: un percorso $x \to y \to z$ con $x \neq y$ e $y \neq z$ non può tornare in $x$, perché un
    ordine largo non ha andate e ritorno tra elementi distinti (antisimmetria).
\end{spiegazione}

Nel seguito lavoreremo quasi sempre con ordini larghi: un ordine stretto si può sempre ricavare togliendo l'uguaglianza.

% ──────────────────────────────────────────────────────────────
\section{Insiemi ordinati}\label{sec:insiemi-ordinati}

\dfn[insieme-ordinato]{Insieme ordinato, ordine indotto}{
    Sia $s$ un insieme e sia $\rho = (s, s, g) \in \mathrm{OL}(s)$.
    \begin{itemize}
        \item La coppia $(s, \rho)$ si dice \textbf{insieme ordinato}.
        \item Se $t \subseteq s$, la relazione $\rho_t = (t, t, g \cap (t \times t))$ si dice \textbf{ordine indotto} da $\rho$ su $t$, e
              $(t, \rho_t)$ si dice \textbf{sottoinsieme ordinato} di $(s, \rho)$. Con abuso di notazione scriveremo spesso $\rho$ invece di $\rho_t$.
    \end{itemize}
    Per un ordine generico useremo spesso il simbolo $\leq$, e scriveremo $x < y$ per $x \leq y \wedge x \neq y$ (cioè $x \leq^\wedge y$).
}

\begin{concetto}
\begin{Proposizione}{L'ordine indotto è un ordine}{ordine-indotto}
    \begin{prerequisiti}
        \dip{Definizione}{def:insieme-ordinato}
        \dip{Definizione}{def:tipi-relazioni}
    \end{prerequisiti}
    Se $(s, \rho)$ è un insieme ordinato e $t \subseteq s$, allora $\rho_t \in \mathrm{OL}(t)$. Inoltre, per $x, y \in t$, $x \rho_t y \leftrightarrow x \rho y$.
\end{Proposizione}
\begin{dimostrazione}
    Per $x, y \in t$ si ha $(x, y) \in t \times t$, quindi $(x, y) \in g \cap (t \times t) \leftrightarrow (x, y) \in g$: questa è la seconda
    affermazione. Le tre proprietà di $\rho_t$ sono allora quelle di $\rho$, applicate a elementi che stanno in $t$: ad esempio se
    $x, y, z \in t$, $x \rho_t y$ e $y \rho_t z$, allora $x \rho y$ e $y \rho z$, quindi $x \rho z$, cioè $x \rho_t z$.
\end{dimostrazione}
\end{concetto}

\begin{spiegazione}
    Restringere un ordine a una parte vuol dire \textbf{cancellare dal grafo gli elementi che non stanno in $t$}, con tutte le
    frecce che li toccano. Le frecce rimaste collegano elementi di $t$ esattamente come prima, quindi riflessività, antisimmetria
    e transitività restano vere. Ad esempio la divisibilità su $\mbN$ induce un ordine su $\{1, 2, 3, 4, 6, 12\}$.
\end{spiegazione}

\dfn[confrontabili]{Elementi confrontabili, ordine totale}{
    Sia $(s, \rho)$ un insieme ordinato.
    \begin{itemize}
        \item Due elementi $x, y \in s$ si dicono \textbf{confrontabili} se $x \rho y \vee y \rho x$.
        \item Se ogni coppia di elementi di $s$ è confrontabile, $\rho$ si dice \textbf{ordine totale} e $(s, \rho)$ si dice
              \textbf{insieme totalmente ordinato} (o \textbf{catena}).
    \end{itemize}
}

\dfn[minimo-minimale]{Minimo, massimo, minimali, massimali}{
    Sia $(s, \rho)$ un insieme ordinato e sia $m \in s$. Diciamo che $m$ è
    \begin{itemize}
        \item \textbf{minimo} di $s$ se $\forall x \in s\,(m \rho x)$: sta sotto a tutti; si scrive $m = \min(s)$;
        \item \textbf{massimo} di $s$ se $\forall x \in s\,(x \rho m)$: sta sopra a tutti; si scrive $m = \max(s)$;
        \item \textbf{minimale} se $\forall x \in s\,(x \rho m \rightarrow x = m)$: non c'è nessuno strettamente sotto di lui;
        \item \textbf{massimale} se $\forall x \in s\,(m \rho x \rightarrow x = m)$: non c'è nessuno strettamente sopra di lui.
    \end{itemize}
}

\begin{concetto}
\begin{Proposizione}{Altre forme della minimalità}{forme-minimale}
    \begin{prerequisiti}
        \dip{Definizione}{def:minimo-minimale}
        \dip{Definizione}{def:confrontabili}
        \dip{Teorema}{teo:negazione-quantificatori}
        \dip{Proposizione}{prop:negimp}
    \end{prerequisiti}
    Sia $(s, \rho)$ un insieme ordinato e $m \in s$. Sono equivalenti:
    \begin{enumerate}
        \item $m$ è minimale;
        \item $\neg\exists y \in s\,(y \rho m \wedge y \neq m)$: non esiste un elemento strettamente più piccolo di $m$;
        \item $\forall x \in s\,\big((x \rho m \vee m \rho x) \rightarrow m \rho x\big)$: ogni elemento confrontabile con $m$ sta sopra $m$.
    \end{enumerate}
    Per i massimali valgono le affermazioni con $\rho$ scambiata con la duale.
\end{Proposizione}
\begin{dimostrazione}
    \begin{caso}{(1) $\leftrightarrow$ (2)}
        Negando (2) si ottiene $\exists y\,(y \rho m \wedge y \neq m)$, che è la negazione di $\forall y\,(y \rho m \rightarrow y = m)$
        (Teorema~\ref{teo:negazione-quantificatori} e Proposizione~\ref{prop:negimp}). Quindi (2) è equivalente a (1).
    \end{caso}
    \begin{caso}{(1) $\rightarrow$ (3)}
        Sia $x$ confrontabile con $m$. Se $m \rho x$ abbiamo finito; se $x \rho m$, per la minimalità $x = m$ e quindi $m \rho x$ per riflessività.
    \end{caso}
    \begin{caso}{(3) $\rightarrow$ (1)}
        Sia $x \rho m$. Allora $x$ è confrontabile con $m$, quindi per (3) $m \rho x$; per l'antisimmetria $x = m$.
    \end{caso}
\end{dimostrazione}
\end{concetto}

\begin{spiegazione}
    «Minimale» non vuol dire «sotto a tutti», ma «\textbf{sotto a tutti quelli con cui si confronta}»: un elemento minimale può
    avere accanto elementi con cui non è confrontabile. Nell'insieme ordinato qui sotto (disegnato con la convenzione dei
    diagrammi di Hasse della Sezione~\ref{sec:hasse}: «più in basso» = «più piccolo») $a$ e $b$ sono entrambi minimali, ma
    nessuno dei due è minimo: non sono confrontabili tra loro.
    \begin{center}
    \begin{tikzpicture}[hasse]
        \node[hv evid] (a) at (0,0) {$a$}; \node[hv evid] (b) at (1.6,0) {$b$};
        \node[hv] (c) at (0.8,1.1) {$c$}; \node[hv] (d) at (0.8,2.2) {$d$};
        \draw[hlato] (a) -- (c) -- (d); \draw[hlato] (b) -- (c);
        \node[font=\small, align=left, anchor=west] at (2.6,1.1)
            {$a \leq c \leq d$, \ $b \leq c$;\\ $a$ e $b$ non sono confrontabili:\\ sono minimali, ma non c'è minimo.\\ $d$ è il massimo.};
    \end{tikzpicture}
    \end{center}
\end{spiegazione}

\begin{concetto}
\begin{Teorema}{Unicità di minimo e massimo}{unicita-minimo}
    \begin{prerequisiti}
        \dip{Definizione}{def:minimo-minimale}
        \dip{Definizione}{def:tipi-relazioni}
    \end{prerequisiti}
    Sia $(s, \rho)$ un insieme ordinato. Il minimo di $s$, se esiste, è unico. Lo stesso vale per il massimo.
\end{Teorema}
\begin{dimostrazione}
    Siano $m_1$ e $m_2$ minimi di $s$. Poiché $m_1$ è minimo, $m_1 \rho m_2$; poiché $m_2$ è minimo, $m_2 \rho m_1$. Per l'antisimmetria
    $m_1 = m_2$. Per il massimo la dimostrazione è identica, con $x \rho m$ al posto di $m \rho x$.
\end{dimostrazione}
\end{concetto}

\begin{spiegazione}
    Due minimi dovrebbero stare \textbf{ciascuno sotto l'altro}; in un ordine questo succede solo se sono lo stesso elemento.
    È lo stesso ragionamento che fa dire che due numeri con $x \leq y$ e $y \leq x$ sono uguali.
\end{spiegazione}

\begin{concetto}
\begin{Teorema}{Il minimo è l'unico minimale}{minimo-minimale}
    \begin{prerequisiti}
        \dip{Definizione}{def:minimo-minimale}
        \dip{Definizione}{def:tipi-relazioni}
    \end{prerequisiti}
    Sia $(s, \rho)$ un insieme ordinato con minimo $m$. Allora $m$ è minimale ed è l'unico elemento minimale di $s$.
    Allo stesso modo il massimo, se esiste, è l'unico elemento massimale.
\end{Teorema}
\begin{dimostrazione}
    \begin{caso}{$m$ è minimale}
        Se $x \rho m$, poiché $m$ è il minimo vale anche $m \rho x$, quindi $x = m$ per l'antisimmetria.
    \end{caso}
    \begin{caso}{Ogni minimale è $m$}
        Sia $x$ minimale. Poiché $m$ è il minimo, $m \rho x$; poiché $x$ è minimale, da $m \rho x$ segue $m = x$.
    \end{caso}
    Per il massimo si ripete il ragionamento con la relazione duale.
\end{dimostrazione}
\end{concetto}

\begin{spiegazione}
    Se c'è un elemento sotto a tutti, nessun altro può essere «senza nessuno sotto»: ha sotto almeno il minimo. Attenzione al
    viceversa: avere un \textbf{unico} minimale non basta in generale per avere un minimo (lo vedremo nel
    Capitolo~\ref{cap:insiemi-ordinati}: negli insiemi finiti basta, in quelli infiniti no).
\end{spiegazione}

\dfn[ben-ordinato]{Insieme ben ordinato}{
    Un insieme ordinato $(s, \rho)$ si dice \textbf{ben ordinato} (e $\rho$ si dice \textbf{buon ordine}) se ogni parte non vuota di
    $s$ ha minimo (rispetto all'ordine indotto).
}

\begin{concetto}
\begin{Teorema}{Ogni buon ordine è totale}{ben-ordinato-totale}
    \begin{prerequisiti}
        \dip{Definizione}{def:ben-ordinato}
        \dip{Definizione}{def:confrontabili}
        \dip{Proposizione}{prop:ordine-indotto}
    \end{prerequisiti}
    Ogni insieme ben ordinato è totalmente ordinato.
\end{Teorema}
\begin{dimostrazione}
    Sia $(s, \rho)$ ben ordinato e siano $x, y \in s$. La parte $\{x, y\}$ non è vuota, quindi ha un minimo $m$, che è $x$ oppure $y$.
    Se $m = x$, allora $x \rho y$; se $m = y$, allora $y \rho x$. In ogni caso $x$ e $y$ sono confrontabili.
\end{dimostrazione}
\end{concetto}

\begin{spiegazione}
    Il buon ordine chiede il minimo di \textbf{ogni} parte non vuota, in particolare delle parti con due elementi: e il minimo di
    $\{x, y\}$ è proprio quello dei due che sta sotto l'altro. Il viceversa è falso: $(\mbZ, \leq)$ è totalmente ordinato ma la
    parte $\{\dots, -2, -1, 0\}$ non ha minimo.
\end{spiegazione}

\begin{esempio}{Insiemi ordinati notevoli}
    \begin{itemize}
        \item $(\mbN, \leq)$ è ben ordinato: ogni parte non vuota di $\mbN$ ha un elemento più piccolo di tutti.
        \item $(\mbZ, \leq)$ è totalmente ordinato ma non ben ordinato, e non ha né minimo né massimo.
        \item Per ogni insieme $s$, $(\mcP(s), \subseteq)$ ha minimo $\varnothing$ e massimo $s$. Se $s$ ha due elementi distinti $x$ e $y$,
              non è totalmente ordinato: $\{x\}$ e $\{y\}$ non sono confrontabili.
        \item $(\mbN, \mid)$ ha minimo $1$ (che divide tutti) e massimo $0$ (che è diviso da tutti: $0 = 0 \cdot n$), ma non è
              totalmente ordinato ($2 \nmid 3$ e $3 \nmid 2$).
        \item In $(\mbN \setminus \{1\}, \mid)$ non c'è minimo; gli elementi minimali sono i numeri che non hanno divisori
              diversi da $1$ e da sé stessi, cioè i numeri primi.
    \end{itemize}
\end{esempio}

\begin{concetto}
\begin{Proposizione}{Quando l'insieme delle parti è totalmente ordinato}{parti-totale}
    \begin{prerequisiti}
        \dip{Definizione}{def:confrontabili}
        \dip{Definizione}{def:inclusione}
        \dip{Proposizione}{prop:contrapposizione}
    \end{prerequisiti}
    Sia $s$ un insieme. $(\mcP(s), \subseteq)$ è totalmente ordinato se e solo se $s = \varnothing$ oppure $s$ è un singleton.
\end{Proposizione}
\begin{dimostrazione}
    \begin{caso}{$(\leftarrow)$}
        $\mcP(\varnothing) = \{\varnothing\}$ ha un solo elemento, che è confrontabile con sé stesso. $\mcP(\{x\}) = \{\varnothing, \{x\}\}$ e
        $\varnothing \subseteq \{x\}$.
    \end{caso}
    \begin{caso}{$(\rightarrow)$, per contrapposizione}
        Se $s$ non è vuoto né un singleton, contiene due elementi distinti $x \neq y$. Allora $\{x\} \not\subseteq \{y\}$ (perché
        $x \notin \{y\}$) e $\{y\} \not\subseteq \{x\}$: $\{x\}$ e $\{y\}$ non sono confrontabili.
    \end{caso}
\end{dimostrazione}
\end{concetto}

\begin{spiegazione}
    Basta che $s$ abbia due elementi perché nascano due «rami» paralleli: i diagrammi di Hasse qui sotto (Sezione~\ref{sec:hasse})
    mostrano le parti di $\varnothing$, di $\{a\}$ e di $\{a, b\}$; solo nel terzo compaiono elementi non confrontabili.
    \begin{center}
    \begin{tikzpicture}[hasse]
        \node[hv] (v0) at (0,0) {$\varnothing$};
        \node[font=\small] at (0,-0.75) {$\mcP(\varnothing)$};
        \begin{scope}[xshift=2.6cm]
            \node[hv] (v1) at (0,0) {$\varnothing$}; \node[hv] (a1) at (0,1.1) {$\{a\}$};
            \draw[hlato] (v1) -- (a1);
            \node[font=\small] at (0,-0.75) {$\mcP(\{a\})$};
        \end{scope}
        \begin{scope}[xshift=5.8cm]
            \node[hv] (v2) at (0,0) {$\varnothing$}; \node[hv evid] (a2) at (-0.9,1.1) {$\{a\}$}; \node[hv evid] (b2) at (0.9,1.1) {$\{b\}$};
            \node[hv] (ab2) at (0,2.2) {$\{a, b\}$};
            \draw[hlato] (v2) -- (a2) -- (ab2); \draw[hlato] (v2) -- (b2) -- (ab2);
            \node[font=\small] at (0,-0.75) {$\mcP(\{a, b\})$};
        \end{scope}
    \end{tikzpicture}
    \end{center}
\end{spiegazione}

% ──────────────────────────────────────────────────────────────
\section{Applicazioni crescenti e isomorfismi}\label{sec:crescenti}

Come per le strutture algebriche, anche tra insiemi ordinati ci interessano le funzioni che ne «rispettano la forma».

\dfn[crescente]{Applicazione crescente, isomorfismo di insiemi ordinati}{
    Siano $(s, \rho)$ e $(s', \rho')$ insiemi ordinati e $f \colon s \to s'$.
    \begin{itemize}
        \item $f$ si dice \textbf{crescente} se $\forall x, y \in s\,\big(x \rho y \rightarrow f(x) \rho' f(y)\big)$.
        \item $f$ si dice \textbf{isomorfismo} (di insiemi ordinati) se è biettiva e $\forall x, y \in s\,\big(x \rho y \leftrightarrow f(x) \rho' f(y)\big)$.
              In questo caso $(s, \rho)$ e $(s', \rho')$ si dicono \textbf{isomorfi}.
    \end{itemize}
}

\begin{esempio}{Isomorfismi}
    \begin{itemize}
        \item $n \in \mbN \mapsto 2n \in \{2n \mid n \in \mbN\}$ è un isomorfismo tra $(\mbN, \leq)$ e l'insieme dei pari con $\leq$:
              è biettiva e $m \leq n \leftrightarrow 2m \leq 2n$.
        \item $n \in \mbN \mapsto 2^n \in \{2^n \mid n \in \mbN\}$ è un isomorfismo tra $(\mbN, \leq)$ e le potenze di $2$ ordinate per
              divisibilità: $m \leq n \leftrightarrow 2^m \mid 2^n$.
    \end{itemize}
\end{esempio}

Ogni isomorfismo è crescente. Al contrario di quanto succede per le strutture algebriche, però, una funzione biettiva e
crescente non è detto che sia un isomorfismo.

\begin{concetto}
\begin{Proposizione}{Una biezione crescente che non è un isomorfismo}{biettiva-crescente}
    \begin{prerequisiti}
        \dip{Definizione}{def:crescente}
        \dip{Definizione}{def:funzioni-notevoli}
    \end{prerequisiti}
    Sia $s = \mbN \setminus \{0\}$. L'identità $\mathrm{id}_s$ è una funzione biettiva e crescente da $(s, \mid)$ a $(s, \leq)$, ma non è
    un isomorfismo.
\end{Proposizione}
\begin{dimostrazione}
    L'identità è biettiva. È crescente: se $m \mid n$ con $n \geq 1$, allora $n = km$ con $k \geq 1$, quindi $m \leq km = n$.
    Non è un isomorfismo: $2 \leq 3$, ma $2 \nmid 3$; cioè $\mathrm{id}(2) \leq \mathrm{id}(3)$ non implica $2 \mid 3$.
\end{dimostrazione}
\end{concetto}

\begin{spiegazione}
    Un isomorfismo deve portare le frecce in frecce \textbf{e} le «non frecce» in «non frecce». Una funzione crescente garantisce
    solo la prima metà: può \textbf{aggiungere} confronti che prima non c'erano. Ecco i diagrammi di Hasse di $\{1, 2, 3, 4, 6\}$
    con la divisibilità e con $\leq$: l'identità porta ogni segmento di sinistra in un percorso verso l'alto a destra, ma a destra
    ci sono confronti nuovi, come $2 \leq 3$ e $4 \leq 6$.
    \begin{center}
    \begin{tikzpicture}[hasse]
        \node[hv] (1) at (0,0) {$1$}; \node[hv] (2) at (-0.8,1.1) {$2$}; \node[hv] (3) at (0.8,1.1) {$3$};
        \node[hv] (4) at (-1.6,2.2) {$4$}; \node[hv] (6) at (0,2.2) {$6$};
        \draw[hlato] (1) -- (2) -- (4); \draw[hlato] (2) -- (6); \draw[hlato] (1) -- (3) -- (6);
        \node[font=\small] at (-0.4,-0.75) {$(\{1, 2, 3, 4, 6\}, \mid)$};
        \draw[-{Stealth}, line width=1pt, primary] (1.7,1.1) -- node[above, font=\small] {$\mathrm{id}$} (3.1,1.1);
        \begin{scope}[xshift=4.3cm, yshift=-0.2cm]
            \node[hv] (a) at (0,0) {$1$}; \node[hv] (b) at (0,0.6) {$2$}; \node[hv] (c) at (0,1.2) {$3$};
            \node[hv] (d) at (0,1.8) {$4$}; \node[hv] (e) at (0,2.4) {$6$};
            \draw[hlato] (a) -- (b); \draw[hlato evid] (b) -- (c); \draw[hlato] (c) -- (d); \draw[hlato evid] (d) -- (e);
            \node[font=\small] at (0,-0.55) {$(\{1, 2, 3, 4, 6\}, \leq)$};
        \end{scope}
    \end{tikzpicture}
    \end{center}
\end{spiegazione}

\begin{concetto}
\begin{Proposizione}{Inversa e composta di isomorfismi}{isomorfismi-ordinati}
    \begin{prerequisiti}
        \dip{Definizione}{def:crescente}
        \dip{Teorema}{teo:biettiva-inversa}
        \dip{Corollario}{cor:comp-biettive}
    \end{prerequisiti}
    L'inversa di un isomorfismo di insiemi ordinati è un isomorfismo, e la composta di due isomorfismi è un isomorfismo.
    Quindi «essere isomorfi» si comporta come una relazione di equivalenza tra insiemi ordinati.
\end{Proposizione}
\begin{dimostrazione}
    \begin{caso}{Inversa}
        Sia $f \colon (s, \rho) \to (s', \rho')$ un isomorfismo; $f^{-1}$ è biettiva. Per $u, v \in s'$, posto $x = f^{-1}(u)$, $y = f^{-1}(v)$:
        $f^{-1}(u) \rho f^{-1}(v) \leftrightarrow x \rho y \leftrightarrow f(x) \rho' f(y) \leftrightarrow u \rho' v$.
    \end{caso}
    \begin{caso}{Composta}
        Se $f \colon (s, \rho) \to (s', \rho')$ e $g \colon (s', \rho') \to (s'', \rho'')$ sono isomorfismi, $g \circ f$ è biettiva
        (Corollario~\ref{cor:comp-biettive}) e $x \rho y \leftrightarrow f(x) \rho' f(y) \leftrightarrow g(f(x)) \rho'' g(f(y))$.
    \end{caso}
    Infine l'identità di $(s, \rho)$ è un isomorfismo di $(s, \rho)$ in sé stesso.
\end{dimostrazione}
\end{concetto}

\begin{spiegazione}
    Un isomorfismo è un \textbf{cambio di nomi} che lascia intatto il disegno: l'inversa rimette i nomi vecchi, e due cambi di nome
    uno dopo l'altro sono ancora un cambio di nome. Per questo due insiemi ordinati isomorfi sono, per tutto ciò che riguarda
    l'ordine, «lo stesso insieme con etichette diverse».
\end{spiegazione}

% ──────────────────────────────────────────────────────────────
\section{Relazione di copertura e diagrammi di Hasse}\label{sec:hasse}

Il grafo con frecce di un ordine ha moltissime frecce: tutti i cappi e tutte le scorciatoie. Per disegnare un insieme ordinato
finito basta molto meno: le frecce tra elementi \textbf{immediatamente} uno sopra l'altro.

\dfn[copertura]{Relazione di copertura, diagramma di Hasse}{
    Sia $(s, \leq)$ un insieme ordinato e siano $x, y \in s$. Si dice che \textbf{$y$ copre $x$} se
    \[ x < y \ \wedge\ \neg\exists z \in s\,(x < z \wedge z < y), \]
    cioè $y$ sta sopra $x$ e non c'è niente in mezzo: $y$ è un \textbf{successore immediato} di $x$. La relazione binaria
    $\gamma$ su $s$ definita da $x \gamma y \leftrightarrow y \text{ copre } x$ si chiama \textbf{relazione di copertura} (o
    \textbf{diagramma di Hasse}) di $(s, \leq)$; è antiriflessiva.

    Il diagramma di Hasse di un insieme ordinato finito si \textbf{disegna} mettendo ogni elemento più in basso degli elementi che lo
    coprono e tracciando un segmento da $x$ a $y$ ogni volta che $y$ copre $x$.
}

\begin{concetto}
Per \textbf{leggere} un diagramma di Hasse: $x \leq y$ se e solo se $x = y$ oppure da $x$ si arriva a $y$ \textbf{salendo} lungo i
segmenti (per gli insiemi finiti lo dimostreremo nel Teorema~\ref{teo:hasse-isomorfismo}). Ad esempio:
\begin{center}
\begin{tikzpicture}[hasse]
    \node[hv] (v) at (0,0) {$\varnothing$}; \node[hv] (e) at (0,1.1) {$\{\varnothing\}$};
    \draw[hlato] (v) -- (e);
    \node[font=\small, align=center] at (0,-0.8) {$(\mcP(\mcP(\varnothing)), \subseteq)$};
    \begin{scope}[xshift=4.2cm]
        \node[hv] (v1) at (0,0) {$\varnothing$};
        \node[hv] (a) at (-1.5,1.1) {$\{a\}$}; \node[hv] (b) at (0,1.1) {$\{b\}$}; \node[hv] (c) at (1.5,1.1) {$\{c\}$};
        \node[hv] (ab) at (-1.5,2.2) {$\{a, b\}$}; \node[hv] (ac) at (0,2.2) {$\{a, c\}$}; \node[hv] (bc) at (1.5,2.2) {$\{b, c\}$};
        \node[hv] (abc) at (0,3.3) {$\{a, b, c\}$};
        \draw[hlato] (v1) -- (a); \draw[hlato] (v1) -- (b); \draw[hlato] (v1) -- (c);
        \draw[hlato] (a) -- (ab); \draw[hlato] (a) -- (ac); \draw[hlato] (b) -- (ab); \draw[hlato] (b) -- (bc);
        \draw[hlato] (c) -- (ac); \draw[hlato] (c) -- (bc);
        \draw[hlato] (ab) -- (abc); \draw[hlato] (ac) -- (abc); \draw[hlato] (bc) -- (abc);
        \node[font=\small] at (0,-0.8) {$(\mcP(\{a, b, c\}), \subseteq)$};
    \end{scope}
    \begin{scope}[xshift=9.4cm]
        \node[hv] (1) at (0,0) {$1$}; \node[hv] (2) at (-0.8,1.1) {$2$}; \node[hv] (3) at (0.8,1.1) {$3$};
        \node[hv] (4) at (-1.6,2.2) {$4$}; \node[hv] (6) at (0,2.2) {$6$}; \node[hv] (12) at (-0.8,3.3) {$12$};
        \draw[hlato] (1) -- (2) -- (4) -- (12); \draw[hlato] (1) -- (3) -- (6) -- (12); \draw[hlato] (2) -- (6);
        \node[font=\small] at (-0.4,-0.8) {divisori di $12$ con $\mid$};
    \end{scope}
    \begin{scope}[xshift=13.2cm]
        \node[hv] (t3) at (0,0) {$3$}; \node[hv] (t6) at (0,0.8) {$6$}; \node[hv] (t23) at (0,1.6) {$23$}; \node[hv] (t1000) at (0,2.4) {$1000$};
        \draw[hlato] (t3) -- (t6) -- (t23) -- (t1000);
        \node[font=\small] at (0,-0.8) {$(\{3, 6, 23, 1000\}, \leq)$};
    \end{scope}
\end{tikzpicture}
\end{center}
Nel secondo diagramma $\{a\} \subseteq \{a, b, c\}$ perché si sale da $\{a\}$ a $\{a, b\}$ e poi a $\{a, b, c\}$; $\{a\}$ e $\{b, c\}$ non sono
confrontabili perché nessuno dei due si raggiunge salendo dall'altro. Un insieme totalmente ordinato finito ha come diagramma
di Hasse una \textbf{catena} verticale, come l'ultimo.
\end{concetto}

\begin{esempio}{Minimi e minimali sul diagramma di Hasse}
    Sul diagramma di Hasse di un insieme ordinato finito:
    \begin{itemize}
        \item un elemento è \textbf{minimale} se non ha segmenti che scendono, \textbf{massimale} se non ha segmenti che salgono;
        \item c'è il \textbf{minimo} se c'è un elemento da cui si raggiungono tutti salendo, il \textbf{massimo} se c'è un elemento
              che si raggiunge salendo da tutti.
    \end{itemize}
    Nei divisori di $12$ il minimo è $1$ e il massimo è $12$. Togliendo $1$ e $12$ restano $\{2, 3, 4, 6\}$: i minimali sono $2$ e $3$,
    i massimali $4$ e $6$, e non ci sono né minimo né massimo.
\end{esempio}

% ──────────────────────────────────────────────────────────────
\section{Minoranti e maggioranti}\label{sec:minoranti}

\dfn[minoranti]{Minoranti e maggioranti}{
    Sia $(s, \rho)$ un insieme ordinato e sia $t \subseteq s$. Un elemento $m \in s$ si dice
    \begin{itemize}
        \item \textbf{minorante} di $t$ se $\forall z \in t\,(m \rho z)$;
        \item \textbf{maggiorante} di $t$ se $\forall z \in t\,(z \rho m)$.
    \end{itemize}
    Indichiamo con $\mathrm{Minor}_{(s, \rho)}(t)$ l'insieme dei minoranti di $t$ e con $\mathrm{Maggior}_{(s, \rho)}(t)$ quello dei
    maggioranti (omettendo il pedice quando non c'è ambiguità).
}

Attenzione: i minoranti e i maggioranti sono elementi di $s$, non necessariamente di $t$. Il minimo di $t$, se esiste, è un
minorante di $t$ che sta in $t$.

\begin{esempio}{Minoranti e maggioranti}
    \begin{itemize}
        \item In $(\mbN, \leq)$ i minoranti di una parte non vuota $t$ sono i numeri minori o uguali a tutti gli elementi di $t$.
        \item In $(\mbN, \mid)$ i minoranti di $t$ sono i \textbf{divisori comuni} di tutti gli elementi di $t$, i maggioranti i
              \textbf{multipli comuni}. Ad esempio $\mathrm{Maggior}(\{2\})$ è l'insieme dei multipli di $2$ (compreso $0$) e
              $\mathrm{Minor}(\{12, 18\}) = \{1, 2, 3, 6\}$.
        \item In $(\mcP(s), \subseteq)$, se $t$ è una parte non vuota di $\mcP(s)$, i minoranti di $t$ sono gli insiemi contenuti in tutti
              gli elementi di $t$, cioè le parti di $\bigcap t$; i maggioranti sono gli insiemi che contengono tutti gli elementi di $t$,
              cioè le parti di $s$ che contengono $\bigcup t$.
        \item Per la parte vuota: ogni elemento di $s$ è sia minorante sia maggiorante di $\varnothing$ (la condizione $\forall z \in \varnothing$
              è vera a vuoto), quindi $\mathrm{Minor}(\varnothing) = \mathrm{Maggior}(\varnothing) = s$.
    \end{itemize}
    Sul diagramma dei divisori di $12$, con $t = \{4, 6\}$ (in arancione): i minoranti (in verde) sono gli elementi da cui si
    raggiungono \textbf{sia} $4$ \textbf{sia} $6$ salendo, i maggioranti quelli che si raggiungono salendo da entrambi.
    \begin{center}
    \begin{tikzpicture}[hasse]
        \node[hv evid2] (1) at (0,0) {$1$}; \node[hv evid2] (2) at (-0.8,1.1) {$2$}; \node[hv] (3) at (0.8,1.1) {$3$};
        \node[hv evid] (4) at (-1.6,2.2) {$4$}; \node[hv evid] (6) at (0,2.2) {$6$}; \node[hv] (12) at (-0.8,3.3) {$12$};
        \draw[hlato] (1) -- (2) -- (4) -- (12); \draw[hlato] (1) -- (3) -- (6) -- (12); \draw[hlato] (2) -- (6);
        \node[font=\small, align=left, anchor=west] at (1.6,1.6)
            {$\mathrm{Minor}(\{4, 6\}) = \{1, 2\}$\\[1mm] $\mathrm{Maggior}(\{4, 6\}) = \{12\}$\\[1mm]
             ($3$ non è un minorante: $3 \nmid 4$)};
    \end{tikzpicture}
    \end{center}
\end{esempio}
